\documentclass[12pt]{article}

% ============================================
% PACKAGES
% ============================================

% Page layout
\usepackage[margin=1in]{geometry}

% Font and encoding
\usepackage[T1]{fontenc}
\usepackage[utf8]{inputenc}
\usepackage{mathptmx}  % Times-like font

% Math
\usepackage{amsmath}
\usepackage{amssymb}

% Graphics and figures
\usepackage{graphicx}
\usepackage{float}
\usepackage{subcaption}

% Tables
\usepackage{booktabs}
\usepackage{tabularx}
\usepackage{multirow}
\usepackage{array}

% Citations and bibliography
\usepackage{natbib}
\bibliographystyle{apa} 

% Hyperlinks (load last among most packages)
\usepackage{hyperref}
\hypersetup{
    colorlinks=true,
    linkcolor=black,
    citecolor=black,
    urlcolor=blue,
    breaklinks=true
}

% Line spacing
\usepackage{setspace}
\onehalfspacing

% For appendix
\usepackage{appendix}

% For landscape pages (if needed for wide tables)
\usepackage{pdflscape}

% ============================================
% CUSTOM COMMANDS
% ============================================

% For easier table notes
\newcommand{\tablenote}[1]{\par\smallskip\footnotesize #1}

% ============================================
% TITLE PAGE SETUP
% ============================================

\title{\Large\textbf{The Impact of Lottery Incentives:\\ 
Increasing Response Rates Can Decrease Survey Accuracy}}

\author{
    Jerold Pearson\thanks{Department of Political Science, University Name. Email: author1@university.edu} 
    \and 
    Roger E. Levine\thanks{Department of Economics, Other University. Email: author2@other.edu}
    \and
    Jon A. Krosnick\thanks{Department of Sociology, Third University. Email: author3@third.edu}
    \and
    Resty Fufunan\thanks{School of Public Policy, Fourth University. Email: author4@fourth.edu}
    \and
    Josearmando Torres\thanks{Department of Government, Fifth University. Email: author5@fifth.edu}
}

\date{\today}

% ============================================
% DOCUMENT BEGINS
% ============================================

\begin{document}

% Title
\maketitle

% ============================================
% ABSTRACT
% ============================================

\begin{abstract}
\noindent
Your abstract goes here. This should be a concise summary of your paper, typically 150-250 words. It should briefly describe the research question, methods, key findings, and implications. The abstract is often the most-read part of your paper, so make it count. Include the main contribution and why readers should care about your findings.
\end{abstract}

\thispagestyle{empty}  % No page number on title page
\newpage
\setcounter{page}{1}

% ============================================
% INTRODUCTION
% ============================================

\section{Introduction}

Your introduction goes here. Start by motivating the research question and establishing why it matters. What puzzle or gap in the literature are you addressing?

Cite relevant work like this \citep{Author2020} or as \citet{Author2020} showed. You can cite multiple sources \citep{Author2020, OtherAuthor2019}.

The introduction should:
\begin{itemize}
    \item Establish the importance of the topic
    \item Identify the specific research question
    \item Preview your approach and findings
    \item Outline the paper's structure
\end{itemize}

% ============================================
% THEORY / LITERATURE REVIEW
% ============================================

\section{Theoretical Framework}

Present your theoretical framework here. What mechanisms do you propose? What hypotheses do you derive?

\subsection{First Mechanism}

Describe the first mechanism or theoretical argument.

\subsection{Second Mechanism}

Describe additional mechanisms as needed.

% ============================================
% DATA AND METHODS
% ============================================

\section{Data and Methods}

Describe your data sources and methodological approach.

\subsection{Data Sources}

Our primary data source is [describe your data]. Table~\ref{tab:descriptive} presents descriptive statistics.

% Example table
\begin{table}[H]
\centering
\caption{Descriptive Statistics}
\label{tab:descriptive}
\begin{tabular}{lccccc}
\toprule
Variable & Mean & SD & Min & Max & N \\
\midrule
Variable 1 & 0.50 & 0.25 & 0 & 1 & 1,000 \\
Variable 2 & 3.45 & 1.23 & 1 & 7 & 1,000 \\
Variable 3 & 45.2 & 15.8 & 18 & 85 & 1,000 \\
\bottomrule
\end{tabular}
\tablenote{\textit{Note:} Data from [source]. Sample restricted to [criteria].}
\end{table}

\subsection{Survey Measures}

If you use survey questions, present them clearly:

\begin{quote}
On [topic], some people argue that [position A]. Others argue that [position B]. Still, others are somewhere in between. Where would you place yourself on this scale, or haven't you thought much about this?
\end{quote}

Respondents place themselves on a 1 to 7 scale, where 1 indicates [position A] and 7 indicates [position B].

\subsection{Empirical Strategy}

We employ [method] to estimate [what you're estimating]:

\begin{equation}
Y_i = \alpha + \beta X_i + \mathbf{Z}_i'\gamma + \varepsilon_i
\label{eq:main}
\end{equation}

\noindent where $Y_i$ is the outcome for individual $i$, $X_i$ is the key independent variable, and $\mathbf{Z}_i$ is a vector of control variables.

% ============================================
% RESULTS
% ============================================

\section{Results}

Present your main findings here. Figure~\ref{fig:main} shows [description].

% Example figure placeholder
\begin{figure}[H]
\centering
% \includegraphics[width=0.8\textwidth]{your_figure.pdf}
\fbox{\parbox{0.8\textwidth}{\centering\vspace{3cm}[Your Figure Here]\vspace{3cm}}}
\caption{Main Results Figure}
\label{fig:main}
\tablenote{\textit{Note:} Description of what the figure shows. Points show [what]. Lines show [confidence intervals, etc.].}
\end{figure}

Table~\ref{tab:regression} presents regression results.

% Example regression table
\begin{table}[H]
\centering
\caption{Regression Results}
\label{tab:regression}
\begin{tabular}{lccc}
\toprule
 & Model 1 & Model 2 & Model 3 \\
\midrule
Key Variable & 0.127$^{*}$ & 0.135$^{*}$ & 0.142$^{*}$ \\
 & (0.012) & (0.013) & (0.014) \\
Control 1 & & 0.007 & 0.005 \\
 & & (0.012) & (0.011) \\
Control 2 & & & $-$0.029 \\
 & & & (0.111) \\
Intercept & $-$0.463$^{*}$ & $-$0.455$^{*}$ & $-$0.448$^{*}$ \\
 & (0.221) & (0.225) & (0.228) \\
\midrule
R$^2$ & 0.354 & 0.356 & 0.358 \\
N & 267 & 267 & 267 \\
\bottomrule
\end{tabular}
\tablenote{$^{*}p < 0.05$. Standard errors in parentheses.}
\end{table}

\subsection{Robustness Checks}

Describe any robustness checks or alternative specifications.

% ============================================
% DISCUSSION / CONCLUSION
% ============================================

\section{Discussion and Conclusion}

Summarize your findings and discuss their implications.

What are the main takeaways? How do your findings contribute to existing literature? What are the limitations? What questions remain for future research?

% ============================================
% REFERENCES
% ============================================

\newpage
\singlespacing

% Option 1: Use BibTeX (recommended)
% \bibliography{references}

% Option 2: Manual bibliography (for template demonstration)
\begin{thebibliography}{99}

\bibitem[Author(2020)]{Author2020}
Author, A. (2020). 
\newblock Article title here.
\newblock \textit{Journal Name}, 45(3), 123--145.

\bibitem[Other Author(2019)]{OtherAuthor2019}
Other Author, B. \& Third Author, C. (2019). 
\newblock Book title here.
\newblock Publisher Name.

\end{thebibliography}

% ============================================
% APPENDIX
% ============================================

\newpage
\onehalfspacing
\appendix
\section*{Appendix}
\addcontentsline{toc}{section}{Appendix}

\renewcommand{\thesection}{\Alph{section}}
\renewcommand{\thetable}{\Alph{section}\arabic{table}}
\renewcommand{\thefigure}{\Alph{section}\arabic{figure}}
\setcounter{section}{0}
\setcounter{table}{0}
\setcounter{figure}{0}

\section{Additional Tables and Figures}

Present supplementary analyses here.

% Example appendix table
\begin{table}[H]
\centering
\caption{Appendix Table: Additional Results}
\label{tab:appendix1}
\begin{tabular}{lcc}
\toprule
 & Full Sample & Subsample \\
\midrule
Variable & 0.50 & 0.55 \\
 & (0.10) & (0.12) \\
\midrule
N & 1,000 & 500 \\
\bottomrule
\end{tabular}
\end{table}

\section{Robustness Analyses}

Additional robustness checks can go here.

\section{Variable Definitions}

\begin{itemize}
    \item \textbf{Variable 1:} Definition and coding.
    \item \textbf{Variable 2:} Definition and coding.
\end{itemize}

\end{document}